\documentclass[10pt,a4paper]{amsart}
\usepackage[margin=19mm]{geometry}
\usepackage{amsmath,amssymb,mathtools}
\usepackage[scr=rsfs,cal=euler]{mathalfa}
\usepackage{xfrac}
\usepackage[unicode,hidelinks,bookmarks=false]{hyperref}
\newcommand{\ie}{i.e.\ }
\newcommand{\cf}{cf.\ }
\newcommand{\resp}{respectively\ }
\newcommand{\Subsection}{\subsection}
\providecommand{\nobreakdash}{\nobreak\hbox{-}}
\setlength{\emergencystretch}{2em}
\multlinegap0pt
\usepackage{bm}
\newtheorem{theorem}{Theorem}
\title{Approximation of a Multivariate Function of Bounded Variation from its Scattered Data}
\author{Rajesh Dachiraju}
\date{Source: arXiv:2011.11258v6 (CC BY 4.0)}
\begin{document}
\maketitle

\noindent\emph{Source and licence.} Approximation of a Multivariate Function of Bounded Variation from its Scattered Data. Version arXiv:2011.11258v6 (CC BY 4.0), distributed by arXiv under the Creative Commons Attribution 4.0 International licence, http://creativecommons.org/licenses/by/4.0/, as recorded in the arXiv licence metadata for this exact version and shown on the abstract page https://arxiv.org/abs/2011.11258v6. Full licence notice: \texttt{licenses/arxiv-2011.11258-CC-BY-4.0.txt}.\par\smallskip

\noindent\textit{Editorial scope.} Counted unit: the article's theorem eig\_theorem (source0.tex lines 456--458). The article's complete original proof of it is delivered with it (source0.tex lines 459--587, one \texttt{proof} environment of 129 lines). No mathematics is written, repaired or re-derived: the statement and the proof are the article's own lines, in the article's own notation, and no retained line is substituted or re-flowed.  Retained uncounted as the source-local context the proof needs: lines 446--448.  Nothing was omitted from any retained span, so the package carries no omission record.  No retained line contains a citation, so the wrapper carries no bibliography and no bibliography entry.  No retained line contains a figure environment, an image-inclusion command or drawing code, so no image is delivered and the package declares no asset dependency.  Editorial formatting only, in addition: an AMS \texttt{amsart} preamble that restates the article's own declarations and macros used by the retained lines, namely the environments \texttt{theorem} on separate counters, together with the standard packages those lines need.  The retained environments print in this document as Theorem 1 in order of appearance, whereas the article prints the same items with the numbering of its own full document.  The complete native source of the article as distributed in the arXiv e-print for this exact version is bundled unchanged as \texttt{sources/104975/source0.tex}.
\begin{equation}\label{eig_G}
\rho_l(G_{\lambda}) = a_{l0} + \frac{a_{l1}}{\lambda} +\frac{a_{l2}}{\lambda^2}+\ldots.
\end{equation}

\begin{theorem}\label{eig_theorem}
 $$\mbox{for    }\lambda>1\mbox{,          } a_{l1} > 0 \mbox{,         } l= 2,3\ldots n.$$ 
\end{theorem}

\begin{proof}
Let $\boldsymbol{v}_l(\lambda) = [v_{l1}(\lambda),v_{l2}(\lambda),v_{l3}(\lambda),\ldots v_{ln}(\lambda)]^T$ be a normalized eigenvector of the matrix $G_{\lambda}$ corresponding to the eigenvalue $\rho_l(G_{\lambda})$ (normalized meaning $\|\boldsymbol{v}_l(\lambda)\|_2=1$). Now, we estimate the eigenvalue $\rho_l(G_{\lambda})$. Thus, by definition of eigenvalue 
\begin{equation}\label{eigest}
\begin{aligned}
    \rho_l(G_{\lambda}) &= \boldsymbol{v}_l^T(\lambda)G_{\lambda}\boldsymbol{v}_l(\lambda)\\
    &= \sum\limits_{i=1}^n\sum\limits_{j=1}^nv_{li}(\lambda)v_{lj}(\lambda)g_{\lambda}(\boldsymbol{p}_i-\boldsymbol{p}_j)\\
    &= \sum\limits_{i=1}^n\sum\limits_{j=1}^n\left\{v_{li}(\lambda)v_{lj}(\lambda)\sum_{\boldsymbol{\eta}\in\mathbb{Z}^m} \frac{1}{1+\lambda\|\boldsymbol{\eta}\|_{2k}^{2k}} \cos{(2\pi\boldsymbol{\eta}\cdot(\boldsymbol{p}_i-\boldsymbol{p}_j))}\right\}.
\end{aligned}
\end{equation} 

we know \begin{equation}\label{dmoiver}
    \cos{(2\pi\boldsymbol{\eta}\cdot(\boldsymbol{p}_i-\boldsymbol{p}_j))} = \frac{1}{2}\left(e^{2\pi i \boldsymbol{\eta}\cdot\boldsymbol{p}_i } e^{-2\pi i \boldsymbol{\eta}\cdot\boldsymbol{p}_j} + e^{-2\pi i \boldsymbol{\eta}\cdot\boldsymbol{p}_i } e^{2\pi i \boldsymbol{\eta}\cdot\boldsymbol{p}_j}\right). 
\end{equation}

Substituting Equation \ref{dmoiver} into Equation \ref{eigest},
\begin{equation}\label{eigenest_cont}
    \begin{aligned}
        \rho_l(G_{\lambda}) &= \boldsymbol{v}_l^T(\lambda)G_{\lambda}\boldsymbol{v}_l(\lambda)\\
        &= \sum\limits_{i=1}^n\sum\limits_{j=1}^n\left\{v_{li}(\lambda)v_{lj}(\lambda)\sum_{\boldsymbol{\eta}\in\mathbb{Z}^m} \frac{1}{1+\lambda\|\boldsymbol{\eta}\|_{2k}^{2k}} \frac{1}{2}\left(e^{2\pi i \boldsymbol{\eta}\cdot\boldsymbol{p}_i } e^{-2\pi i \boldsymbol{\eta}\cdot\boldsymbol{p}_j} + e^{-2\pi i \boldsymbol{\eta}\cdot\boldsymbol{p}_i } e^{2\pi i \boldsymbol{\eta}\cdot\boldsymbol{p}_j}\right)\right\}\\
        %&= \frac{1}{2}\sum\limits_{i=1}^n\sum\limits_{j=1}^n\left\{v_{li}(\lambda)v_{lj}(\lambda)\sum_{\boldsymbol{\eta}\in\mathbb{Z}^m} \frac{1}{1+\lambda\|\boldsymbol{\eta}\|_{2k}^{2k}}e^{2\pi i \boldsymbol{\eta}\cdot\boldsymbol{p}_i } e^{-2\pi i \boldsymbol{\eta}\cdot\boldsymbol{p}_j} + \sum_{\boldsymbol{\eta}\in\mathbb{Z}^m} \frac{1}{1+\lambda\|\boldsymbol{\eta}\|_{2k}^{2k}}e^{-2\pi i \boldsymbol{\eta}\cdot\boldsymbol{p}_i } e^{2\pi i \boldsymbol{\eta}\cdot\boldsymbol{p}_j} \right\}\\
        &= \frac{1}{2}\sum_{\boldsymbol{\eta}\in\mathbb{Z}^m}\left\{ \frac{1}{1+\lambda\|\boldsymbol{\eta}\|_{2k}^{2k}}\sum\limits_{i=1}^n\sum\limits_{j=1}^n v_{li}(\lambda)v_{lj}(\lambda)\left\{e^{2\pi i \boldsymbol{\eta}\cdot\boldsymbol{p}_i } e^{-2\pi i \boldsymbol{\eta}\cdot\boldsymbol{p}_j} + e^{-2\pi i \boldsymbol{\eta}\cdot\boldsymbol{p}_i } e^{2\pi i \boldsymbol{\eta}\cdot\boldsymbol{p}_j}
\right\} \right\}\\
        &= \sum_{\boldsymbol{\eta}\in\mathbb{Z}^m}\left\{ \frac{1}{1+\lambda\|\boldsymbol{\eta}\|_{2k}^{2k}}\left\{\left( \sum\limits_{i=1}^n v_{li}(\lambda)e^{2\pi i \boldsymbol{\eta}\cdot\boldsymbol{p}_i } \right) \left( \sum\limits_{j=1}^n v_{lj}(\lambda)e^{-2\pi i \boldsymbol{\eta}\cdot\boldsymbol{p}_j } \right) \right\} \right\}. 
\end{aligned}
\end{equation}

Denoting $z_l(\lambda,\boldsymbol{\eta}) = \sum\limits_{i=1}^n v_{li}(\lambda) e^{2\pi i \boldsymbol{\eta}\cdot\boldsymbol{p}_i }$, we have
\begin{equation}\label{eigenest_cont2}
\begin{aligned}
    \rho_l(G_{\lambda}) &= \boldsymbol{v}_l^T(\lambda)G_{\lambda}\boldsymbol{v}_l(\lambda)\\
    &= \sum_{\boldsymbol{\eta}\in\mathbb{Z}^m}\left\{ \frac{z_l(\lambda,\boldsymbol{\eta})\bar{z_l}(\lambda,\boldsymbol{\eta})}{1+\lambda\|\boldsymbol{\eta}\|_{2k}^{2k}}\right\}\\
    &= \sum_{\boldsymbol{\eta}\in\mathbb{Z}^m}\left\{ \frac{1}{1+\lambda\|\boldsymbol{\eta}\|_{2k}^{2k}}\left|z_l(\lambda,\boldsymbol{\eta})\right|^2 \right\}.
\end{aligned}
\end{equation}


$\left|z_l(\lambda,\boldsymbol{\eta})\right|^2 \ge 0$ and $z_l(\lambda,\boldsymbol{\eta})$ cannot vanish for all $\eta\in\mathbb{Z}^m$. Hence, we conclude that
\begin{equation}\label{psd}
\begin{aligned}
   \rho_l(G_{\lambda}) &= \boldsymbol{v}_l^T(\lambda)G_{\lambda}\boldsymbol{v}_l(\lambda)\\
        &= \sum_{\boldsymbol{\eta}\in\mathbb{Z}^m}\left\{ \frac{1}{1+\lambda\|\boldsymbol{\eta}\|_{2k}^{2k}}\left|z_l(\lambda,\boldsymbol{\eta})\right|^2 \right\}\\
        &> 0.
\end{aligned}
\end{equation}
This proves that all eigenvalues of the matrix $G_{\lambda}$ are positive; hence, $G_{\lambda}$ is a positive definite matrix. Now, we need to estimate the eigenvalues asymptotically as $\lambda\to\infty$. Let us consider Equation \ref{psd} and split the sum to obtain
\begin{equation}\label{psd2}
\begin{aligned}
   \rho_l(G_{\lambda}) &= \boldsymbol{v}_l^T(\lambda)G_{\lambda}\boldsymbol{v}_l(\lambda)\\
        &= \sum_{\boldsymbol{\eta}\in\mathbb{Z}^m}\left\{ \frac{1}{1+\lambda\|\boldsymbol{\eta}\|_{2k}^{2k}}\left|z_l(\lambda,\boldsymbol{\eta})\right|^2 \right\}\\
            &= \sum_{\boldsymbol{\eta}\in\{0\}^m}\left\{ \frac{1}{1+\lambda\|\boldsymbol{\eta}\|_{2k}^{2k}}\left|z_l(\lambda,\boldsymbol{\eta})\right|^2 \right\} + \sum_{\boldsymbol{\eta}\in\mathbb{Z}^m\setminus\{0\}^m}\left\{ \frac{1}{1+\lambda\|\boldsymbol{\eta}\|_{2k}^{2k}}\left|z_l(\lambda,\boldsymbol{\eta})\right|^2 \right\}\\
            &= \left|z_l(\lambda,\boldsymbol{0})\right|^2 + \sum_{\boldsymbol{\eta}\in\mathbb{Z}^m\setminus\{0\}^m}\left\{ \frac{1}{1+\lambda\|\boldsymbol{\eta}\|_{2k}^{2k}}\left|z_l(\lambda,\boldsymbol{\eta})\right|^2 \right\}\\
        &= \left|\sum\limits_{i=1}^nv_{li}(\lambda)\right| ^2 + \sum_{\boldsymbol{\eta}\in\mathbb{Z}^m\setminus\{0\}^m}\left\{ \frac{1}{1+\lambda\|\boldsymbol{\eta}\|_{2k}^{2k}}\left|z_l(\lambda,\boldsymbol{\eta})\right|^2 \right\}.   
\end{aligned}
\end{equation}
\begin{equation}\label{seq_1}
 \left|z_l(\lambda,\boldsymbol{\eta})\right|^2 \ge 0
 \end{equation}\begin{equation}\label{proof_lambda}
\begin{aligned}
    z_l(\lambda,\boldsymbol{\eta}) &= \sum\limits_{i=1}^n v_{li}(\lambda) e^{2\pi i \boldsymbol{\eta}\cdot\boldsymbol{p}_i }\\
    \implies \left|z_l(\lambda,\boldsymbol{\eta})\right|^2 &=  \left|\sum\limits_{i=1}^n v_{li}(\lambda) e^{2\pi i \boldsymbol{\eta}\cdot\boldsymbol{p}_i } \right|^2\\
    &\le\sum\limits_{i=1}^n \left| v_{li}(\lambda)\right|^2            \sum\limits_{i=1}^n \left|e^{2\pi i \boldsymbol{\eta}\cdot\boldsymbol{p}_i }  \right|^2\\ 
    &\le n\|\boldsymbol{v}_l\|_2^2\\
    &=n.
\end{aligned}
\end{equation}
\begin{equation}\label{seq_3}
\forall \lambda>0, \sup\limits_{\eta\in\mathbb{Z}^m\setminus\{0\}} \left| z_l(\lambda,\boldsymbol{\eta})  \right|^2 > 0.
\end{equation}
Thus, using Equations \ref{seq_1}, \ref{proof_lambda}, \ref{seq_3} \begin{equation}\label{eq_label}
\mbox{as   } \lambda\to\infty, \sup\limits_{\eta\in\mathbb{Z}^m\setminus\{0\}} \left| z_l(\lambda,\boldsymbol{\eta})  \right|^2 = \Theta(1)
\end{equation}
\begin{equation}\label{eq_proof_eig}
\begin{aligned}
\sum_{\boldsymbol{\eta}\in\mathbb{Z}^m\setminus\{0\}^m}\left\{ \frac{\lambda}{1+\lambda\|\boldsymbol{\eta}\|_{2k}^{2k}}\left|z_l(\lambda,\boldsymbol{\eta})\right|^2 \right\} &\le \lambda \left(\sup\limits_{\eta\in\mathbb{Z}^m\setminus\{0\}} \left| z_l(\lambda,\boldsymbol{\eta})  \right|^2\right) \sum_{\boldsymbol{\eta}\in\mathbb{Z}^m\setminus\{0\}^m}\frac{1}{1+\lambda\|\boldsymbol{\eta}\|_{2k}^{2k}} \\
&\le \left(\sup\limits_{\eta\in\mathbb{Z}^m\setminus\{0\}} \left| z_l(\lambda,\boldsymbol{\eta})  \right|^2\right) \sum_{\boldsymbol{\eta}\in\mathbb{Z}^m\setminus\{0\}^m}\frac{1}{\|\boldsymbol{\eta}\|_{2k}^{2k}}, \\
&\le K \sup\limits_{\eta\in\mathbb{Z}^m\setminus\{0\}} \left| z_l(\lambda,\boldsymbol{\eta})  \right|^2,
\end{aligned}
\end{equation} where $K = \sum\limits_{\boldsymbol{\eta}\in\mathbb{Z}^m\setminus\{0\}^m}\frac{1}{\|\boldsymbol{\eta}\|_{2k}^{2k}} \mbox{            }$ is a finite positive constant as $2k>1$ as it was already assumed in the beginning that $k>\frac{m}{2}$.

Using Equations \ref{eq_label} and \ref{eq_proof_eig},
\begin{equation}\label{forward_eq}
\lambda\sum_{\boldsymbol{\eta}\in\mathbb{Z}^m\setminus\{0\}^m}\left\{ \frac{1}{1+\lambda\|\boldsymbol{\eta}\|_{2k}^{2k}}\left|z_l(\lambda,\boldsymbol{\eta})\right|^2 \right\} < K_1 \mbox{   (a positive constant)}.
\end{equation}

As $\left|z_l(\lambda,\boldsymbol{\eta})\right|^2$ is always positive as long as $\lambda>0$ and does not vanish for all $\boldsymbol{\eta}\in\mathbb{Z}^m\setminus\{0\}$, there exists an $\boldsymbol{\eta}_0$, for which it is positive. Let \begin{equation}\label{delta}
\left|z_l(\lambda,\boldsymbol{\eta}_0)\right|^2 = \delta_0 > 0.
\end{equation}
Hence, \begin{equation}\label{reverse}
\begin{aligned}
\lambda\sum_{\boldsymbol{\eta}\in\mathbb{Z}^m\setminus\{0\}^m}\left\{ \frac{1}{1+\lambda\|\boldsymbol{\eta}\|_{2k}^{2k}}\left|z_l(\lambda,\boldsymbol{\eta})\right|^2 \right\} &\ge  \frac{\lambda\delta_0}{1+\lambda\|\boldsymbol{\eta}_0\|_{2k}^{2k}}\\
&> \frac{\delta_0}{1+\|\boldsymbol{\eta}_0\|_{2k}^{2k}}\mbox{               ( as  }\lambda >1 \mbox{ as stated in the theorem)       }, \\
&= K_0 > 0
\end{aligned}
\end{equation}
Combining Equations \ref{forward_eq} and \ref{reverse},
%we already know that $\sum_{\boldsymbol{\eta}\in\mathbb{Z}^m\setminus\{0\}^m}\left\{ \frac{1}{1+\lambda\|\boldsymbol{\eta}\|_{2k}^{2k}}\left|z_l(\lambda,\boldsymbol{\eta})\right|^2 \right\}$ is always positive; hence, there exists a $K_0>0$ such that $$\sum_{\boldsymbol{\eta}\in\mathbb{Z}^m\setminus\{0\}^m}\left\{ \frac{1} {1+\lambda\|\boldsymbol{\eta}\|_{2k}^{2k}}\left|z_l(\lambda,\boldsymbol{\eta})\right|^2 \right\}>K_0$$
\begin{equation}
\label{pp}
\begin{aligned}
\frac{K_0}{\lambda} < \sum_{\boldsymbol{\eta}\in\mathbb{Z}^m\setminus\{0\}^m}\left\{ \frac{1}{1+\lambda\|\boldsymbol{\eta}\|_{2k}^{2k}}\left|z_l(\lambda,\boldsymbol{\eta})\right|^2 \right\} &< \frac{K_1}{\lambda} 
\end{aligned}
\end{equation}
Thus, $$\rho_l(G_{\lambda}) = \left|\sum\limits_{i=1}^nv_{li}(\lambda)\right|^2 +\sum_{\boldsymbol{\eta}\in\mathbb{Z}^m\setminus\{0\}^m}\left\{ \frac{1}{1+\lambda\|\boldsymbol{\eta}\|_{2k}^{2k}}\left|z_l(\lambda,\boldsymbol{\eta})\right|^2 \right\} = \left|\sum\limits_{i=1}^nv_{li}(\lambda)\right|^2 +\Theta(\frac{1}{\lambda}).$$ 

Let \begin{equation}
    \begin{aligned}
        h(\lambda) &= \left|\sum\limits_{i=1}^nv_{li}(\lambda)\right|^2\\\mbox{  and }\\
        u(\lambda) &= \sum_{\boldsymbol{\eta}\in\mathbb{Z}^m\setminus\{0\}^m}\left\{ \frac{1}{1+\lambda\|\boldsymbol{\eta}\|_{2k}^{2k}}\left|z_l(\lambda,\boldsymbol{\eta})\right|^2 \right\}.
        \end{aligned}
\end{equation}
Hence, \begin{equation}\label{sum}
   \rho_l(G_{\lambda}) = h(\lambda) + u(\lambda).
\end{equation}

As both $u(\lambda)$ and $h(\lambda)$ are positive and that for $l>1$ $\rho_l(G_{\lambda}) \to 0$ as $1/\lambda\to 0$, we have \begin{equation}\label{gotozero}
    \begin{aligned}
        \lim\limits_{\frac{1}{\lambda}\to0}h(\lambda) &= 0\\  \mbox{  and  }\\
        \lim\limits_{\frac{1}{\lambda}\to0}u(\lambda) &= 0.
    \end{aligned}
\end{equation}

Let $h(\lambda) = h_a(\lambda) + h_n(\lambda)$, where $h_a(\lambda)$ is an analytic function of $\frac{1}{\lambda}$ in any neighborhood of $1/\lambda =\epsilon = 0$ and $h_n(\lambda)$ is not an analytic function of $\frac{1}{\lambda}$ in any neighborhood of $1/\lambda =\epsilon = 0$. Similarly, let $u(\lambda) = u_a(\lambda) + u_n(\lambda)$. From Equation \ref{pp}, $u(\lambda) = \Theta(1/\lambda)$. Hence, $u_a(\lambda)$ is of the form $u_a(\lambda) = \frac{w_1}{\lambda}+\frac{w_2}{\lambda^2}+\cdots$, where $w_1>0$ and $u_n(\lambda) = o(1/\lambda)$. As $\rho_l(G_{\lambda})$ is an analytic function of $1/\lambda$, by using Equation \ref{sum}, we have $h_n(\lambda) = - u_n(\lambda)$, and there is $\left|h_n(\lambda)\right| = o(1/\lambda)$. Using this and Equation \ref{gotozero}, we can say that $h_a(\lambda)$ is of the form $\frac{x_1}{\lambda}+\frac{x_2}{\lambda^2}+\cdots$. Thus, as $h(\lambda)$ is always positive and $|h_n(\lambda)| = o(1/\lambda)$, we have that $\lambda h_a(\lambda)$ is positive for sufficiently large $\lambda$. Thus, $\lambda(\frac{x_1}{\lambda}+\frac{x_2}{\lambda^2}+\ldots)$ is positive for sufficiently large $\lambda$, meaning $x_1\ge 0$. Thus, using Equation \ref{sum} and that $h_n(\lambda) = -u_n(\lambda)$, we have \begin{equation}\label{new_sum}
    \rho_l(G_{\lambda}) = \left(\frac{x_1}{\lambda}+\frac{x_2}{\lambda^2}+\cdots\right) + \left(\frac{w_1}{\lambda}+\frac{w_2}{\lambda^2}+\cdots\right).
\end{equation} As we have already shown that $w_1>0$ and $x_1\ge 0$, comparing Equations \ref{new_sum} and \ref{eig_G}, for $l>1$, as $a_{l0} = 0$, we have $a_{l1} = w_1+x_1 > 0$. Hence, it is proved.


%As $\|\boldsymbol{v}_l(\lambda)\|_2 = 1$, we have $\sum\limits_{i=1}^nv_{li}(\lambda) = O(1) $ Hence for $\lambda>1$,  $a_{l1} > 0 \mbox{,               } l= 1,2,3\ldots n$

\end{proof}

\end{document}
